% !TEX program = xelatex
% eprint 风格中文预印本 —— v3（按 2026-10-14 审阅意见修订）
% 核心修订：max_beta=1754 归因、完整成本口径、撤回 fatigue 定性、
%           逐参数集清单与容差、DTRU 因果收紧、文献更正
\documentclass[11pt,a4paper,zihao=-4]{ctexart}
\usepackage[top=2.5cm,bottom=2.5cm,left=2.6cm,right=2.6cm]{geometry}
\usepackage{amsmath,amssymb}
\usepackage{booktabs}
\usepackage{xcolor}
\usepackage{enumitem}
\usepackage[colorlinks=true,linkcolor=blue!50!black,citecolor=blue!50!black,urlcolor=blue!50!black]{hyperref}
\usepackage{caption}
\captionsetup{font=small,labelfont=bf}

\setlength{\parindent}{2em}
\setlength{\parskip}{0.25em}

\newcommand{\code}[1]{\texttt{#1}}

\title{\bfseries 如何审计格方案的安全强度声称：来自 NICCS 四卷 KEM\\
复现的三个教训\\[0.8em]
\large\mdseries Auditing Claimed Security Strengths of Lattice KEMs:\\
Three Lessons from Reproducing Four NICCS Submissions}
\author{Jy L}
\date{2026 年 10 月 15 日}

\begin{document}
\maketitle

\begin{abstract}
\noindent\textbf{摘要。}设计文档声称的安全强度与独立复现之间的差距
常被当作单一事实使用，但差距的成因至少有三类，且修正路径完全不同：
复现侧的\textbf{输入还原}误差、估算工具的\textbf{搜索域限制}、以及
\textbf{成本口径}的不完整。本文以 NICCS 四卷候选 KEM（FLIT、ZEN、
NEV、DTRU）推荐的 14 套参数为对象，用各方案引用的工具
（\code{lattice-estimator} 于其指定 commit、Kyber
\code{security-estimates} 脚本）完成一次以差距为中心的审计，
并给出三个量化教训。\textbf{教训一（输入还原）}：四类输入误读
——分布记号歧义、参数表结构误读、分布族近似、$\sigma$ 侧选错——
分别造成 $13\sim26$、$\ge 170$、$13\sim29$、$13\sim41$ bits 的
伪差距；我们首轮复现报告的 $-18\sim-44$ bits「系统性差距」
主要是这些误读的叠加。\textbf{教训二（搜索域）}：
\code{lattice-estimator} 默认搜索上限为
$\beta \le \texttt{max\_beta}{=}1754$（源码注释：对应至多 512-bit
安全等级）；14 套参数中 6 套的声称值高于该等级，默认配置下
返回 \code{rop}$=\infty$，易被误读为「声称不可行」。仅在进程内
将上限提至 2200（不改源码），Flit-512（$\beta{=}1885$，$\log_2
\code{rop}{=}550.42$，与声称 550.4 一致）、ZEN-512 NTRU
（1802，526.18 vs 526.2）、NEV-R3（1999，583.7 vs 583）即由
优化器直接复现。\textbf{教训三（完整成本）}：混合攻击的成本
是完整 \code{rop}（含猜测/FFT/重复项），不是 $0.292\beta$ 的
块大小基础值——ZEN-512 RLWE 面按完整成本与声称仅差
$\approx 4.1$ bits，而非我们此前误算的 21.5 bits。
修正全部口径后：14 套参数中 13 套与公共工具一致（容差内；
其中 5 套 $\beta$ 精确一致），唯一实质残差是 DTRU-648 的
15.7 bits 未解释差异（将输入 $\sigma^2$ 从 1 改为 2 可得
163.8$\approx$声称 164，但数值反解不具唯一性，不作因果结论）。
文中所有比较按方案自身口径逐面进行，不做跨方案排序；
文末给出可复核的配置与产物清单。
\\[0.6em]
\noindent\textbf{Abstract.} A gap between a scheme's claimed security
strength and an independent reproduction can arise from at least
three causes with entirely different remedies: input-restoration
errors on the reproducer's side, search-domain limits of the
estimation tool, and incomplete cost accounting. Reproducing the 14
recommended parameter sets of four NICCS KEM submissions (FLIT,
ZEN, NEV, DTRU) with their own referenced tools, we quantify three
lessons. \emph{Lesson 1 (inputs):} four classes of misreading
---notation ambiguity, table-structure misreading, distribution
approximation, wrong $\sigma$ side---produce spurious gaps of
$13$--$26$, $\ge 170$, $13$--$29$ and $13$--$41$ bits; the
$-18$ to $-44$-bit ``systematic gap'' of our own first pass is
mostly their superposition. \emph{Lesson 2 (search domain):} the
lattice-estimator caps its search at \code{max\_beta}${}=1754$
(source comment: security levels of at most 512 bits); six of the
14 sets claim above that level and return \code{rop}$=\infty$
under default settings, easily misread as infeasibility of the
claims. Raising the cap to 2200 in-process (no source change)
reproduces Flit-512 ($\beta{=}1885$, $\log_2$\code{rop}${}=550.42$
vs claimed 550.4), ZEN-512 NTRU (1802, 526.18 vs 526.2) and NEV-R3
(1999, 583.7 vs 583) directly via the optimizer. \emph{Lesson 3
(full cost):} a hybrid attack costs its full \code{rop} (guessing,
FFT, repetitions), not the $0.292\beta$ block-size floor; ZEN-512's
RLWE face differs from its claim by only $\approx 4.1$ bits under
full cost, not the 21.5 bits we first computed. After all
corrections, 13 of 14 sets agree with the public tools within
stated tolerances (five with exact $\beta$); the sole substantive
residual is a 15.7-bit unexplained discrepancy for DTRU-648
(input $\sigma^2{=}2$ yields $163.8\approx$ the claim 164, but
numerical inversion is non-unique and no causal claim is made).
All comparisons follow each scheme's own caliber per attack face;
no cross-scheme ranking is attempted. A reproducible artifact
list is included.\\[0.6em]
\noindent\textbf{关键词。}格密码；安全强度估算；可复现性；
审计方法；NTRU；RLWE；NICCS\\[0.2em]
\noindent\textbf{Keywords.} lattice cryptography; concrete
security; reproducibility; audit methodology; NTRU; RLWE; NICCS
\end{abstract}

\section{引言}

NICCS 候选 KEM 的安全强度声称值由外部工具生成
（\code{lattice-estimator} 指定 commit、Kyber
\code{security-estimates} 脚本），估算代码不随包提供，独立复现
因此成为评审的关键环节。但「复现有差距」这个观察本身几乎不含
信息：差距可能来自复现侧输入还原、工具搜索域、成本口径或
方案侧真实偏差。把成因搞错会导向错误行动——把工具限制当成
方案缺陷，或把口径差当成攻击面。

本文报告一次\textbf{以差距为中心}的完整审计，覆盖 FLIT、ZEN、
NEV、DTRU 各自\textbf{推荐}参数集共 14 套（DTRU-Light/Prime、
ZEN-light/backup、NEV 的 C/D 系列不在范围内；选择规则是各文档
标记为推荐/主集者）。我们的贡献是三个可量化的教训与一套可
复核的流程，而非新攻击或对任何方案安全性的否定：
\begin{enumerate}[leftmargin=2em,itemsep=0.1em]
\item 输入还原误差的 bit 级量化（\S\ref{sec:l1}），含我们自己
      首轮复现的完整错误清单——这些误读在无随包脚本的条件下
      极易发生；
\item 工具搜索域限制的识别与影响（\S\ref{sec:l2}）：
      \code{max\_beta} 上限使 $>512$-bit 声称默认返回
      \code{rop}$=\infty$；提升上限后三处声称由优化器直接复现；
\item 完整成本口径的必要性（\S\ref{sec:l3}）：以 ZEN-512 为例，
      $0.292\beta$ 与完整 \code{rop} 相差 17.4 bits，直接决定
      「差距 21.5」还是「差距 4.1」；
\item 修正后的逐套清单（\S\ref{sec:inventory}）：13/14 一致
      （5 套 $\beta$ 精确），唯一实质残差 DTRU-648；
\item 可复核产物与配置（\S\ref{sec:artifact}）。
\end{enumerate}

本文不做跨方案安全等级横向比较（四方案口径互异，
\S\ref{sec:l3}）；所有复现值均为估算层面，不构成实际攻击，
也不提供任何「安全下界」——搜索截断与启发式模型预测
不等价于安全性证明。

\section{审计设置}

\textbf{参数与约定。}参数取自四卷文档的参数表；分布记号有歧义
时以\textbf{参考实现}裁决：FLIT 的 $T$ 依其 \code{params.h} 与
Algorithm 2 判为每符号概率（方差 $2T$），其规范散文的「占比」
表述与实现矛盾（两种读法均测试）；ZEN 的 $S_a$ 定义明确
（$\Pr[\pm1]{=}a$ 各取）；NEV 的 1409 是 R 系列模数 $q$；
DTRU 脚本输入秘密侧 $\sigma{=}\sqrt{k_r/2}$。和分布
（如 $B_1{+}S_{3/32}$）用方差匹配的高斯代理——\textbf{代理
等价性是近似}：固定重量三值、独立三值与高斯在方差之外的性质
不同，涉及猜测的攻击（dual\_hybrid）对其敏感
（\S\ref{sec:l3}）。

\textbf{工具。}lattice-estimator 于 FLIT/ZEN 指定 commit
（\code{6019056}，2026-04-28；\code{3e48ef4}，2026-06-27）与
当前 master（\code{53da598}，2026-08-19），Sage 10.9；
DTRU 用其引用的 Kyber 脚本。攻击面按各方案自述配置：RLWE 跑
\code{primal\_usvp} 与 \code{dual\_hybrid}，NTRU 跑
\code{primal\_bdd} 与「NTRU $=$ LWE$(b{=}0)$」建模（NEV 自述
方法）；成本模型 ADPS16（CoreSVP）为主、MATZOV 对照。
Kyber-512 校准：$\beta{=}406$，与已知一致~\cite{newhope,kyber}。

\section{教训一：输入还原误差的量化}\label{sec:l1}

我们首轮复现曾报告 FLIT/ZEN/DTRU 系统性偏低 $18\sim44$ bits、
NEV 出现 $\pm 170\sim 300$ bits 异常。逐项勘误后，这些「差距」
的主体是复现侧输入误读的叠加（表~\ref{tab:sens}）。

\begin{table}[htbp]\centering\small
\caption{四类输入误读的量化影响（首轮 $\to$ 修正后，bits）}\label{tab:sens}
\begin{tabular}{@{}llll@{}}
\toprule
误读类 & 案例 & 首轮 $\to$ 后 & 影响 \\
\midrule
记号歧义 & FLIT $T$ 按散文读（方差 $T$） & $105.7 \to 118.8$；$241.2 \to 267.2$ & $-13.1/-26.0$ \\
表结构误读 & NEV 把 $q{=}1409$ 读作 $n$ & $309.2/280.6 \to 135.8/583.7$ & $+173/-303$ \\
分布族近似 & ZEN $S_a$ 方差取 $a$（应 $2a$） & $99.3 \to 127.6$ 等 & $-13 \sim -29$ \\
$\sigma$ 侧选错 & DTRU 输入 $\sqrt{1/2}$（应 $\sqrt{k_r/2}$） & $135.1 \to 148.3$ 等 & $-13 \sim -41$ \\
\bottomrule
\end{tabular}
\end{table}

其中 FLIT 的记号歧义源自其文档自身：规范散文（「nonzero 系数
占比」）、Algorithm 2 伪代码与参考实现 \code{params.h}、安全表
三方不一致——后三方一致（每符号 $\sigma$），散文是错的。该
不一致已向征集方通报为文档问题；对复现者而言，按散文读会
系统性低估其安全性。

\section{教训二：搜索域上限}\label{sec:l2}

lattice-estimator 的 \code{conf.py} 在三个使用版本中均设置
\begin{quote}\small
\code{max\_beta = 1754 \# Upper bound on blocksize, for security\\
levels of at most 512 bits.}
\end{quote}
\code{primal\_usvp} 与 dual 系列的优化范围由它限定。14 套参数
中 6 套的声称安全等级高于 512 bits（Flit-512、ZEN-512、
NEV-R3、DTRU-1536/2048 等），其最优 $\beta$（1802--1999）超出
默认搜索域，于是默认配置返回 \code{rop}$=\infty$。我们首轮把
其解读为「声称不可行/不可复现」，并猜测了错误的机制（网格
漏搜「可行薄片」）——该猜测已被否定并撤回。

正确的诊断与修复很简单：\textbf{在进程内提升上限}（不改源码）：
\begin{verbatim}
import estimator.lwe_primal as LP, estimator.lwe_dual as LD
LP.max_beta_global = 2200; LD.max_beta_global = 2200
\end{verbatim}
提升后，优化器直接给出（默认 ADPS16，$d{=}n{+}m{+}1$）：

\begin{table}[htbp]\centering\small
\caption{默认上限 vs 上限 2200（优化器直接输出）}\label{tab:cap}
\begin{tabular}{@{}llrrll@{}}
\toprule
实例（建模） & 声称 & \multicolumn{2}{c}{上限 2200} & 一致性 & 默认 \\
 &  & $\beta$ & $\log_2$\code{rop} & （差） & 上限 \\
\midrule
Flit-512（LWE，固定重量代理 var $7/8$） & 550.4 & 1885 & 550.420 & $+0.02$ & $\infty$ \\
ZEN-512 NTRU（LWE$(b{=}0)$，var $3/16$） & 526.2 & 1802 & 526.184 & $-0.02$ & $\infty$ \\
NEV-R3（固定重量代理） & 583 & 1999 & 583.708 & $+0.7$ & $\infty$ \\
NEV-R3（均匀三值 DG 代理） & 583 & 1999 & 583.708 & $+0.7$ & $\infty$ \\
\bottomrule
\end{tabular}
\end{table}

四组在 \code{53da598} 与 \code{3e48ef4} 上结果相同，Flit-512
另在 \code{6019056} 复核。要点有三：
(1) \code{rop}$=\infty$ 在此的语义是\textbf{搜索域未覆盖}，
不是「实例不可解」或「声称错误」；
(2) 该上限是源码注释明示的设计选择（512-bit 等级），但
\textbf{失败语义未文档化}，跨过它的声称值在默认配置下不可复现；
(3) 我们此前把 NTRU 模块对 $n{=}2048$ 恒返回 $\beta{=}1755$
（$=$ 上限$+1$）解释为「overstretched fatigue 饱和」——该解释
\textbf{不成立}：按工具自身判据
$\texttt{ntru\_fatigue\_lb}(2048){=}416{,}463 \gg q$，这些实例
并非 overstretched；恒定 $\beta$ 来自上限截断。此前的 fatigue
表述一并撤回。

作为诊断手段，「把声称点 $(m,\beta)$ 代入闭式成功条件求值」
仍有效（我们以此先于上限实验定位了问题），但它只是验证
「存在满足启发式条件的攻击点」，不能升级为全局最优性或
方案安全性证明。

\section{教训三：完整成本口径}\label{sec:l3}

$0.292\beta$ 是块大小的基础成本（CoreSVP 口径），primal 攻击
在 ADPS16 下 $\log_2$\code{rop} $= 0.292\beta$ 恰好成立；
但 \code{dual\_hybrid} 的完整 \code{rop} 还含猜测、FFT 与重复
项。ZEN 的头条值是其 estimator 输出的完整成本，与之比较必须
用完整 \code{rop}：

\begin{table}[htbp]\centering\small
\caption{ZEN-512 RLWE：$0.292\beta$ vs 完整 \code{rop}
（\code{3e48ef4}，ADPS16）}\label{tab:fullcost}
\begin{tabular}{@{}lrrrl@{}}
\toprule
口径 & $\beta$ & 值 & 与声称 519.9 之差 & 备注 \\
\midrule
$0.292\beta$（错误口径） & 1707 & 498.444 & $-21.5$ & 已撤回 \\
完整 $\log_2$\code{rop} & 1707 & 515.817 & $\mathbf{-4.1}$ &
$m{=}1197$，重复 $2^{16.9}$，$\zeta{=}336$ \\
含压缩噪声（完整） & 1855 & 588.068 & $+68.2$ & 与文档「忽略压缩」不符 \\
\bottomrule
\end{tabular}
\end{table}

即 ZEN-512 RLWE 面与声称仅差约 $4.1$ bits——在我们代理分布
与配置下的\textbf{未归因小残差}，不构成方向性结论。同理，
ZEN-128/256 的保守偏差按完整成本为 $+14.3/+13.5$ bits
（dual\_hybrid 完整 \code{rop} $143.7/275.4$ vs 声称
$129.4/261.9$）。需强调：dual\_hybrid 依赖秘密分布的支持集与
猜测概率，固定重量三值、独立三值与方差匹配高斯\textbf{并非
自动等价}；「方差相同」不构成精确复现的证明，本表数值均以
所用代理为条件。

\section{修正后的逐套清单}\label{sec:inventory}

「复现成功」的判定规则（预先声明）：在方案\textbf{自身口径}
与攻击面上，公共工具在显式配置（含搜索上限）下的输出与声称
之差 $\le 4$ bits（$\beta$ 差 $\le 1$ 视为取整一致）。
按套（非按攻击面）计：\textbf{13/14 一致，1 套实质残差}。

\begin{table}[htbp]\centering\small
\caption{逐套清单（C 为该套头条口径下的经典 bits；差 $=$ 复现 $-$ 声称）}\label{tab:inv}
\begin{tabular}{@{}llllr@{}}
\toprule
参数集 & 攻击面/口径 & 复现（$\beta$，$\log_2$\code{rop}） & 声称 & 差 \\
\midrule
Flit-128 & RLWE/NTRU primal，$0.292\beta$ & 407/407 & 118.8 & $\le 0.3$ \\
Flit-256 & 同上 & 915/915 & 267.2 & $0.0$ \\
Flit-512 & 同上（上限 2200） & 1885，550.42 & 550.4 & $+0.02$ \\
ZEN-128 & NTRU primal & 437，127.6 & 128.5 & $-0.9$ \\
ZEN-256 & NTRU primal & 880，257.0 & 257.5 & $-0.5$ \\
ZEN-512 & NTRU primal（上限 2200） & 1802，526.18 & 526.2 & $-0.02$ \\
 & RLWE dual\_hybrid 完整 & 1707，515.82 & 519.9 & $-4.1$ \\
NEV-R1 & NTRU/RLWE primal & 464/465 & 135 & $\le 0.8$ \\
NEV-R2 & 同上 & 974/976 & 284（$\beta{=}975$） & $\pm 1\beta$ \\
NEV-R3 & NTRU primal（上限 2200） & 1999，583.71 & 583 & $+0.7$ \\
DTRU-648 & 脚本 primal & 507，148.29 & 164 & $\mathbf{-15.7}$ \\
DTRU-768 & 同上 & 684，200.1 & 198 & $+2.1$ \\
DTRU-1024 & 同上 & 920，269.1 & 271 & $-1.9$ \\
DTRU-1536 & 同上 & 1394，407.7 & 411 & $-3.3$ \\
DTRU-2048 & 同上 & 1931，564.8 & 567 & $-2.2$ \\
\bottomrule
\end{tabular}
\end{table}

精确一致（$\beta$ 或 $\log_2$\code{rop} 到小数位）的 5 套：
Flit-256、Flit-512、ZEN-512（NTRU 面）、NEV-R1、NEV-R3
（$+0.7$ 为 $\beta$ 取整差）；NEV-R2 差 $\pm 1\beta$，此前
「逐位一致」的说法不准确，已更正。ZEN-128/256 的 RLWE 面
保守偏差（$+14.3/+13.5$）不构成问题。

\textbf{唯一实质残差：DTRU-648}。在明确的脚本/输入
（n$=$m$_{\max}{=}648$，$q{=}3457$，$\sigma{=}1$）下差
$-15.7$ bits；将 $\sigma^2$ 改为 2（$\beta{=}560$）得
$163.79\approx$ 声称 164。可支持的结论：\textbf{在所列配置下
存在约 15.7 bits 的未解释差异，改变方差可得接近声称的数值}。
不可支持的结论：设计者实际用了 CBD-4、或「参数表与估算必有一
错」——不同攻击、优化域与成本约定可产生相近数值，反解不具
唯一性。闭合该残差需原始估算输入。

\section{讨论：对征集流程的要求}\label{sec:disc}

三个教训各对应一项可执行要求：
(1) 随包估算脚本（参数表、commit、分布到工具输入的映射），
使输入还原不再是猜谜（教训一）；
(2) 声称值随附\textbf{完整配置}（含搜索上限等非常规项）与
$(m,\beta)$ 点，使 \code{rop}$=\infty$ 类失败可归因（教训二）；
(3) 头条口径显式标注（成本模型、是否完整 \code{rop}、取最小
方式），混合攻击不得以 $0.292\beta$ 摘要（教训三）。
FLIT 的散文/伪代码/实现三方不一致应作为文档质量案例通报。

\textbf{局限。}和分布与均匀三值采用方差匹配代理，涉及猜测的
攻击对其敏感；搜索上限提升后未做穷尽性复核（如 dual 系列在
上限处的行为）；本文为估算审计，未执行任何实际格攻击，全部
数值是模型预测而非安全下界。

\section{结论}

以差距为中心审计 NICCS 四卷 KEM 的 14 套推荐参数后：首轮的
「系统性差距」被分解为输入误读（已量化）、搜索域上限
（已定位与绕过）与成本口径（已改用完整 \code{rop}）三类成因；
修正后 13/14 套与公共工具一致（5 套精确），唯一实质残差
DTRU-648（15.7 bits，未归因）。三个教训连同逐点核验、上限
实验与完整成本重算，构成一套可复用的估算审计流程；
其自身不构成新攻击或安全证明。

\section*{复现配置与产物}\label{sec:artifact}

\addcontentsline{toc}{section}{复现配置与产物}

\noindent 工具与版本：lattice-estimator commits
\code{6019056}（2026-04-28）、\code{3e48ef4}（2026-06-27）、
\code{53da598}（2026-08-19）；Sage 10.9（micromamba）；
\code{pq-crystals/security-estimates}（Kyber 脚本）。
脚本（随文库存档于 \code{code/}）：
\code{run\_estimates\_corrected.py}（基线电池）、
\code{review\_verify.py}（上限实验与完整成本重算）、
\code{gap\_localize.py}（闭式条件核验）、
\code{run\_dtru\_kyber.py}（DTRU 脚本复现）。运行示例：
\begin{verbatim}
micromamba run -n sage python review_verify.py   # 上限+完整成本
python run_dtru_kyber.py                          # DTRU 五套
\end{verbatim}
全部原始输出（JSON 与日志）随文库存档；上限修改仅在进程内
进行，未改 estimator 源码。

\begin{thebibliography}{99}\small

\bibitem{estimator}
M.~R. Albrecht, et al. (contributors).
\newblock The Lattice Estimator: security estimates for LWE/NTRU.
\newblock Software, \url{https://github.com/malb/lattice-estimator},
  commits \code{6019056} (2026-04-28), \code{3e48ef4} (2026-06-27),
  \code{53da598} (2026-08-19).

\bibitem{newhope}
E. Alkim, L. Ducas, T. Pöppelmann, P. Schwabe.
\newblock Post-quantum Key Exchange---A New Hope.
\newblock \emph{USENIX Security Symposium}, 2016. (CoreSVP 模型与
  uSVP-GSA 条件.)

\bibitem{matzov}
MATZOV.
\newblock Revisiting the Security of Lattice-Based Schemes.
\newblock IACR ePrint 2022/939, 2022.

\bibitem{dp23}
L. Ducas, M. Pulles.
\newblock Does the Dual(-Hybrid) Attack Even Work?
\newblock \emph{CRYPTO}, 2023.

\bibitem{kyber-est}
P. Schwabe et al. (pq-crystals).
\newblock security-estimates: scripts for estimating the security
  of Kyber.
\newblock Software,
  \url{https://github.com/pq-crystals/security-estimates}.

\bibitem{kyber}
J. Bos, L. Ducas, E. Kiltz, T. Lepoint, V. Lyubashevsky, J. Schanck,
P. Schwabe, G. Seiler, D. Stehlé.
\newblock CRYSTALS---Kyber: a CCA-secure module-lattice-based KEM.
\newblock \emph{IEEE EuroS\&P}, 2018.

\bibitem{chen15}
Y. Chen.
\newblock Lattice Algorithms and Cryptanalysis of Lattice-Based
  Cryptography.
\newblock PhD thesis, École Normale Supérieure de Lyon, 2015.

\bibitem{kannan83}
R. Kannan.
\newblock Improved Algorithms for Integer Programming and Related
  Lattice Problems.
\newblock \emph{STOC}, 1983.

\bibitem{bg14}
S. Bai, S. D. Galbraith.
\newblock An Improved Compression Technique for Signatures Based
  on Learning with Errors.
\newblock \emph{CT-RSA} (Topics in Cryptology), LNCS 8366,
  pp.~28--47, 2014. (小秘密 LWE 的 Bai--Galbraith 嵌入.)

\bibitem{flit-spec}
Y. Pan, Y. Wu, Q. Lu, C. Cheng, R. Liu, Q. Du.
\newblock FLIT: A Lattice-Based Key Encapsulation Mechanism---Algorithm
  Specification and Supporting Documentation.
\newblock NICCS 提交文档, 2026.

\bibitem{zen-spec}
X. Lu et al.
\newblock ZEN Algorithm Specifications.
\newblock NICCS 提交文档, 2026.

\bibitem{nev-spec}
NEV 设计组.
\newblock 算法文本：NEV.
\newblock NICCS 提交文档, 2026.

\bibitem{dtru-spec}
DTRU 设计组.
\newblock DTRU Algorithm Specifications.
\newblock NICCS 提交文档, 2026.

\bibitem{fatigue}
L. Ducas, M. van Woerden.
\newblock The NTRU Fatigue.
\newblock \emph{ASIACRYPT}, 2022. (overstretched 判据的出处；
  本文实例未触发该机制.)

\end{thebibliography}

\end{document}
